Roads are an important form of infrastructure investment that affect households' amenity values, firms’ production functions, and people’s commuting costs. A great deal of existing research has focused on the effects of new roads\footnote{See \cite{ghani2016highway}, \cite{beenstock2016hedonic}, \cite{levkovich2016effects}, \cite{hoogendoorn2019house}, \cite{theisen2021road}, \cite{gibbons2019new}, \cite{fretz2022highways}. Some papers have also observed adverse effects of new roads such as state-led repression \citep{gonzalez2025} and environmental degradation \citep{asher2020eco}.}, especially in developing nations, providing valuable policy insights and spurring development initiatives like China’s Belt and Road Initiative \citep{huang2016understanding} and India’s Pradhan Mantri Gram Sadak Yojana \citep{asher2020}. However, the literature is constrained by two fundamental challenges. The first is endogenous selection as infrastructure investments are not exogenous events but are deliberate policy choices. New road placement is selected by policymakers and is endogenous to anticipated economic growth, while repavement schedules and decisions are systematically correlated with confounding local factors such as political cycles or unobserved community characteristics that independently drive economic outcomes \citep{duranton2011fundamental, gonzalez2023land}. This identification problem is exacerbated by the second challenge: measurement. Most studies rely on road quality metrics that are inherently subjective, labor-intensive, and generated at low frequency \citep{attoh2013pavement}. Such data are not only prone to measurement error that attenuates estimates, but their lack of scalability limits the geographic and temporal scope of empirical analysis \citep{list2022voltage}.
 
In this paper, we avoid endogeneity issues by studying property taxes that fund the maintenance of existing roads and by utilizing a quasi-experimental setting that arises from analyzing close elections. We establish a novel dataset that allows us to study how changes in local taxes affect local infrastructure maintenance and neighborhood house prices. We match voting data on local road maintenance taxes from the Ohio Secretary of State with financial records of cities from Ohio Auditor of State and property sale prices from CoreLogic. Our data also includes remote sensory data of local roads in Ohio from Google Earth Pro. Using these rich data, we first estimate the drop in the maintenance budget for a local government after a tax cut from a failed renewal of road maintenance levy. We fine-tune an Artificial Intelligence (AI) vision model using more than 53,000 satellite images from \cite{brewer2021} to predict changes in road quality after a road maintenance tax cut. Finally, we use the appropriate Regression Discontinuity (RD) estimator to estimate the effect of cutting local road maintenance taxes on housing prices.
% \footnote{As \cite{cellini2010value} shows, when elections are independent, the ITT estimator is equal to the Treatment on the Treated (TOT) estimator. Recent papers such as \cite{hsu2024dynamic} and \cite{biasi_what_2025} have focused on improving upon methods from \cite{cellini2010value}. However, these papers focus on the effects of new capital projects rather than maintaining existing capital infrastructure. Given the exogeneity of the timing of renewal elections and our focus on existing maintenance, the ITT estimator is more appropriate for our setting.}.
Our focus on road maintenance instead of new development, our approach of leveraging close elections as a quasi-experiment to study the effects of local road maintenance tax cuts, and our use of an AI vision model to measure changes in road quality distinguish our study from the existing literature.
Our results reveal that when a city cuts its renewal road taxes, it experiences an 11\% loss in maintenance funds, its road quality declines by 23\% and its house prices decrease by around \$15,350 (9\%). After establishing the main results, we explore heterogeneity in the effects. We find that the effects on house prices are driven by urban areas rather than rural areas, with urban areas experiencing a 7.6\% decline in house prices while rural areas experience inconsistent changes in house prices. We also observe that higher-priced homes experience larger percentage reductions in value relative to lower-priced homes. This suggests that wealthier households are more sensitive to road quality, as a poor road in front of a high-priced home is more likely to be noticed than a poor road in front of a low-priced home.

These results are consistent with other work on the impact of roads on house prices. For instance, \cite{gonzalez2016paving} finds that paving roads in Acayucan, Mexico, increases property values by 28\%, while \cite{theisen2021road} finds 13\% higher housing prices in towns nearest to a new highway in Norway.
Unlike \cite{beenstock2016hedonic} and \cite{diao2017spatial}, we do not find any anticipation effects but we do find statistically significant effects starting in the fourth year after the vote and continuing at least through the ninth year. This delayed effect highlights how the consequences of road tax cuts are not immediate but gradually accumulate as funding is reduced and roads deteriorate, with these effects to be more pronounced for urban rather than rural areas. We also find an effect on the intensive margin with evidence for dosage-response based on the size of the tax cut. 

{\bf Roadmap.} The rest of the paper is organized as follows. Section \ref{sec:lit_review} reviews the related literature and positions our contribution. Section \ref{sec:data} provides background information and presents the variables used in the study. Section \ref{sec:empirical} outlines the empirical strategy. Section \ref{sec:result} presents the results of the study and shares the relevant robustness checks. Section \ref{sec:mech} examines the mechanisms underlying our findings and Section \ref{sec:conclusion} concludes. 

% Finally, we develop a Dynamic General Equilibrium (DGE) model to understand the mechanisms behind our empirical findings. The model consists of a representative household, a local government, and a competitive housing market. Households value well-maintained roads, but they also dislike paying higher property taxes to fund road maintenance, and formalizes the trade-off for households between local road maintenance taxes and the housing amenities.



% Many studies examine the importance of roads and how they affect residents' daily lives (\Citeauthor{currier2023} \citeyear{currier2023}; \Citeauthor{adukia2020} \citeyear{adukia2020}; \Citeauthor{asher2020} \citeyear{asher2020}). Some papers emphasize that building new roads increases employment and entry of new firms \citep{gibbons2019new}, thereby increasing economic activity. Other papers point out the potential for worse outcomes in the form of increased inequality between the rich and the poor \citep{hettige2006} and an exodus of workers seeking access to larger labor markets \citep{asher2020}.  
 
% In this paper, we highlight the importance of road maintenance by townships, cities and villages (``cities'') in Ohio. We collect data on more than 3,000 elections to renew local road tax levies in Ohio from 1991 to 2021. Using the Dynamic Regression Discontinuity (RD) design, we examine how these exogenous cuts in maintenance spending influence housing prices, focusing on cities where voting results are narrowly decided. Our results reveal that when a city cuts its renewal road taxes, it experiences an 11\% loss in maintenance funds, its road quality declines by 23\% and its house prices decrease by around \$15,350 (9\%), with the effect starting four years after the vote and persisting through several years after the vote. The delayed effect is consistent with the time it takes for roads to deteriorate noticeably. We find that these price effects are driven by urban areas rather than rural areas. We also observe that higher-priced houses experience larger percentage reductions in value than lower-priced houses, indicating a heterogeneous impact across different quantiles. This suggests that higher-income households are more sensitive to road quality than lower-income households. 

% We find that these price effects are delayed but clear in urban areas, where house prices fall by 7.6\% on average, compared to inconsistent changes in values of houses in rural areas. We also observe that higher-priced homes experience larger reductions in value relative to lower-priced homes, indicating a heterogeneous impact across different quantiles of home prices. This suggests that wealthier households are more sensitive to road quality, as a poor road in front of a high-priced home is more likely to be noticed than a poor road in front of a low-priced home.

